\section{Discussion}

\label{sec:discussion}
% ══════════════════════════════════════════════════════════════════════════════

\subsection{Advantages of Graph-Based Modeling}

HabitatNet's spatial graph representation provides three key advantages
over traditional SDMs:

\paragraph{Spatial context.} Message passing aggregates information from
neighboring cells, allowing the model to recognize habitat edges,
ecotones, and the spatial configuration of suitable patches---information
invisible to point-based models.

\paragraph{Learned connectivity.} By learning edge weights from
occurrence data, HabitatNet discovers dispersal barriers and corridors
that may not be apparent from land cover alone. Analysis of learned edge
weights reveals that the model assigns high resistance to agricultural
land (mean conductance 0.12), moderate resistance to secondary forest
(0.47), and low resistance to primary forest (0.91).

\paragraph{Scale-aware predictions.} Multi-hop message passing allows
the model to incorporate information at multiple spatial scales without
explicit multi-scale feature engineering.

\subsection{Limitations}

Several limitations warrant consideration:

\begin{enumerate}
  \item \textbf{Biotic interactions}: HabitatNet models abiotic habitat
    suitability and does not explicitly account for species interactions
    such as competition, predation, or mutualism.
  \item \textbf{Adaptive capacity}: The model does not represent genetic
    adaptation or phenotypic plasticity, potentially overestimating range
    loss for adaptable species.
  \item \textbf{Land-use change}: Our base projections use static land
    cover. Integration with land-use change scenarios (LUH2) is
    straightforward but introduces additional uncertainty.
  \item \textbf{Computational cost}: Graph construction and GNN training
    on continental-scale grids require significant computational
    resources (approximately 48 GPU-hours for a full training run on
    the Neotropical realm).
  \item \textbf{Data bias}: GBIF occurrence records are spatially biased
    toward accessible areas and wealthy nations
    \citep{beck2014spatial}. While spatial thinning and bias correction
    mitigate this, residual effects may persist.
\end{enumerate}

\subsection{Conservation Implications}

Our results have direct implications for conservation planning:

\begin{enumerate}
  \item \textbf{Protected area expansion}: Current protected area
    networks cover only 34\% of projected future habitat, suggesting
    significant gaps in anticipatory conservation planning.
  \item \textbf{Corridor investment}: Targeted corridor restoration
    offers high returns---recovering 72\% of lost connectivity with only
    5\% additional area---supporting investment in landscape-scale
    conservation.
  \item \textbf{Climate refugia}: HabitatNet identifies climate
    refugia---areas of persistent suitability across scenarios---that
    merit highest protection priority.
\end{enumerate}

\subsection{Open Science and Reproducibility}

All code, trained models, and processed datasets are released through the
Zoo Open Science Network under Apache 2.0 and CC-BY-4.0 licenses,
respectively. The framework is available as a Python package
(\texttt{habitatnet}) with documentation and tutorials. Model training is
fully reproducible through provided configuration files and random seeds.


% ══════════════════════════════════════════════════════════════════════════════
